\subsection{AI Systems that Learn to Deceive or Manipulate}
\label{sec:6.3.2}
Deception in a machine does not require an intention to deceive, and the three settings where it shows up have almost nothing in common except the effect.

A generative model produces convincing imitations of a real person because that is what it was built to do; whether the result is a film effect or an impersonation is a fact about the user. In a multi-agent environment where agents compete for a limited resource, an agent optimizing for reward can learn to feign weakness or emit false signals, because in some environments deception is what wins — and no part of the system had to represent deceiving anyone. Adversarial examples run the other way: the system is the one being fooled, by an input crafted to flip its output, and the practical effect of a bypassed check is the same as if the system had been suborned.

The responses differ accordingly. The first is a governance problem about provenance and disclosure. The second is a design choice about what a training environment rewards. The third is the robustness engineering of section~\ref{sec:6.2.2}. Treating the three as one problem called “deceptive AI” is how each of them gets the wrong remedy.
